\documentclass[aspectratio=169,11pt]{beamer}

\usefonttheme{professionalfonts}
\usepackage{fontspec}
\usepackage{amsmath,amssymb}
\usepackage{mathtools}
\usepackage{xcolor}
\usepackage{fvextra}

\setsansfont{Avenir Next}
\setmonofont{Menlo}

\definecolor{accent}{HTML}{315C8C}
\definecolor{textgray}{HTML}{222222}
\definecolor{leanComment}{HTML}{6A9955}
\definecolor{leanKeyword}{HTML}{9B3FAE}
\definecolor{leanDecl}{HTML}{0B6E99}
\definecolor{leanType}{HTML}{267F99}
\definecolor{leanOperator}{HTML}{B35C00}
\definecolor{leanHole}{HTML}{7A7A7A}

\setbeamercolor{normal text}{fg=textgray,bg=white}
\setbeamercolor{frametitle}{fg=accent,bg=white}
\setbeamercolor{title}{fg=accent}
\setbeamercolor{structure}{fg=accent}
\setbeamerfont{title}{size=\Large,series=\bfseries}
\setbeamerfont{frametitle}{size=\large,series=\bfseries}
\setbeamerfont{section title}{size=\Large,series=\bfseries}
\setbeamertemplate{navigation symbols}{}
\setbeamertemplate{footline}{%
  \hfill\usebeamerfont{page number in head/foot}\insertframenumber\hspace{1.5em}\vspace{1em}}
\setbeamertemplate{itemize item}{\raisebox{0.2ex}{\small\(\bullet\)}}
\setbeamertemplate{itemize subitem}{\raisebox{0.2ex}{\scriptsize\(\bullet\)}}

\newcommand{\partslide}[1]{%
  \begin{frame}[plain]
    \centering
    \vfill
    {\usebeamerfont{section title}\color{accent}#1}
    \vfill
  \end{frame}
}
\newcommand{\interp}[1]{[\![#1]\!]}
\newcommand{\lcom}[1]{\textcolor{leanComment}{#1}}
\newcommand{\lkw}[1]{\textcolor{leanKeyword}{\textbf{#1}}}
\newcommand{\ldecl}[1]{\textcolor{leanDecl}{#1}}
\newcommand{\ltype}[1]{\textcolor{leanType}{#1}}
\newcommand{\lop}[1]{\textcolor{leanOperator}{#1}}
\newcommand{\lhole}[1]{\textcolor{leanHole}{#1}}

\title{Formalizing Quantifier Elimination}
\subtitle{A Lean formalization in model theory}
\author{}
\date{}

\begin{document}

\begin{frame}[plain]
  \titlepage
\end{frame}

\partslide{Model Theory}

\begin{frame}{Language}
  \begin{itemize}
    \item A first-order language specifies relation and function symbols.
    \item Each symbol has a fixed arity.
    \item DLO language: one binary relation symbol, interpreted as \(\leq\).
  \end{itemize}
\end{frame}

\begin{frame}{Structure}
  \begin{itemize}
    \item A structure interprets the symbols of a language on a domain.
    \item For the order language, a structure is a type \(M\) with a binary relation.
    \item Example: \((\mathbb{Q}, \leq)\).
  \end{itemize}
\end{frame}

\begin{frame}{Terms}
  \begin{itemize}
    \item Terms are built from variables and function symbols.
    \item In the pure order language, there are no function symbols.
    \item DLO terms are just variables or parameters: \(x,y,a,b\).
  \end{itemize}
\end{frame}

\begin{frame}{Formula}
  \begin{itemize}
    \item Formulas are built from atomic formulas using logical connectives and quantifiers.
    \item Atomic DLO examples: \(x \leq y\), \(x = y\).
    \item Formula example: \(\exists z\,(x < z \wedge z < y)\).
  \end{itemize}
\end{frame}

\begin{frame}{Sentence}
  \begin{itemize}
    \item A sentence is a formula with no free variables.
    \item Example: \(\forall x\,\forall y\,(x \leq y \vee y \leq x)\).
    \item Density: \(\forall x\,\forall y\,(x < y \to \exists z\,(x < z \wedge z < y))\).
  \end{itemize}
\end{frame}

\begin{frame}{Satisfaction}
  \begin{itemize}
    \item \(M \models \varphi[a]\): \(\varphi\) holds in \(M\) under assignment \(a\).
    \item For sentences, no assignment is needed.
    \item Example: \((\mathbb{Q},\leq) \models \exists z\,(0<z \wedge z<1)\).
  \end{itemize}
\end{frame}

\begin{frame}{Theory}
  \begin{itemize}
    \item A theory is a set of sentences.
    \item The theory DLO contains the axioms for:
      \begin{itemize}
        \item linear orders,
        \item density,
        \item no endpoints.
      \end{itemize}
  \end{itemize}
\end{frame}

\begin{frame}{Model}
  \begin{itemize}
    \item \(M \models T\) means \(M\) satisfies every sentence in \(T\).
    \item \((\mathbb{Q},\leq)\) is a model of DLO.
    \item \((\mathbb{Z},\leq)\) is not: it is not dense.
  \end{itemize}
\end{frame}

\partslide{Type Theory}

\begin{frame}{Judgements}
  \begin{itemize}
    \item \(a : A\) is a judgement: \(a\) has type \(A\).
    \item Every expression has a type.
    \item Definitional equality is written \(a \equiv b\).
  \end{itemize}
\end{frame}

\begin{frame}{Propositions as Types}
  \begin{itemize}
    \item A proposition is viewed as a type of evidence.
    \item A proof of \(P\) is a term \(p : P\).
    \item To prove a theorem is to construct an object of the corresponding type.
  \end{itemize}
\end{frame}

\begin{frame}{Constructive Reasoning}
  \begin{itemize}
    \item A proof of \(\exists x, P(x)\) gives a witness and evidence.
    \item A proof of \(P \vee Q\) chooses one side and proves it.
    \item The law of excluded middle, \(P \vee \neg P\), is not built in constructively.
    \item Lean allows classical reasoning when explicitly invoked.
  \end{itemize}
\end{frame}

\begin{frame}{Curry--Howard Isomorphism}
  \[
  \begin{aligned}
  \interp{P \Rightarrow Q} &\equiv \interp{P} \to \interp{Q} \\
  \interp{P \wedge Q} &\equiv \interp{P} \times \interp{Q} \\
  \interp{\mathrm{True}} &\equiv 1 \\
  \interp{P \vee Q} &\equiv \interp{P} + \interp{Q} \\
  \interp{\mathrm{False}} &\equiv 0 \\
  \interp{\forall x : A.\,P} &\equiv \Pi x : A.\,\interp{P} \\
  \interp{\exists x : A.\,P} &\equiv \Sigma x : A.\,\interp{P}
  \end{aligned}
  \]
\end{frame}

\partslide{Lean}

\begin{frame}[fragile]
\begin{Verbatim}[fontsize=\scriptsize,commandchars=\\\|\@]
\lcom|/-- A formula is equivalent to a quantifier-free formula over a theory. -/@
\lkw|def@ \ldecl|IsQFEquivalent@ (T : L.Theory) {α : \ltype|Type*@} (φ : L.Formula α) : \ltype|Prop@ :=
  \lop|∃@ ψ : L.Formula α, ψ.IsQF \lop|∧@ φ \lop|⇔[T]@ ψ

\lcom|/-- A theory has quantifier elimination if every formula is equivalent,@
\lcom|over the theory, to a quantifier-free formula. -/@
\lkw|def@ \ldecl|HasQuantifierElimination@ (T : L.Theory) : \ltype|Prop@ :=
  \lop|∀@ {α : \ltype|Type@} (φ : L.Formula α), T.IsQFEquivalent φ
\end{Verbatim}
\end{frame}

\begin{frame}[fragile]
\begin{Verbatim}[fontsize=\scriptsize,commandchars=\\\|\@]
\lcom|/-- The theory of dense linear orders without endpoints@
\lcom|has quantifier elimination. -/@
\lkw|theorem@ \ldecl|dlo_hasQuantifierElimination@ :
    Language.order.dlo.HasQuantifierElimination := \lkw|by@
  \lhole|...@
\end{Verbatim}
\end{frame}

\end{document}
